\section{Alderman Fixed Effects and Index Validation}
\label{sec:appendix_alderman_fe}

This appendix provides additional detail on how the aldermanic strictness index is constructed and validated.

\subsection{Construction of the Strictness Index}

Alderman strictness is measured using permit processing times from the high-discretion permit sample described in Section~\ref{subsec:permits_sample}: new construction, major renovation, demolition, porch construction, and reinstated permits.
The score uses applications through December 2022 and residualizes processing times to remove the influence of ward characteristics.

\paragraph{Stage 1: Residualization.}
I regress permit-level log processing time on ward geographic fundamentals (distance to the CBD, distance to the lake, CTA station density), demographic controls (homeownership rate, population, median household income, racial composition), lagged ward-month permit volume, and fixed effects for month, permit type, and review type.

\paragraph{Stage 2: Alderman Fixed Effects.}
I aggregate the Stage 1 residuals to ward-month mean residuals and regress those ward-month averages on alderman fixed effects $\alpha_{a(w,m)}$ and the same lagged permit-volume control, weighting observations by ward-month permit counts.
Here $a(w,m)$ denotes the alderman serving ward $w$ in month $m$.
The alderman fixed effect captures the systematic component of processing time attributable to the alderman.

\paragraph{Empirical Bayes Shrinkage.}
Raw alderman fixed effects are shrunk toward the grand mean using an empirical Bayes procedure: $\hat{\theta}_j^{\text{EB}} = B_j \cdot \hat{\theta}_j$, where $B_j = \hat{\tau}^2/(\hat{\tau}^2 + \hat{\sigma}_j^2)$ is the shrinkage factor.

\paragraph{Index Standardization.}
The final strictness index is standardized to have mean zero and unit standard deviation across aldermen.
Higher values indicate stricter aldermen.

Table~\ref{tab:stage1_residualization} reports the baseline Stage 1 residualization regression.
The dependent variable is permit-level log processing time.
The specification controls for ward fundamentals and lag permit volume while partialing out month, permit-type, and review-type fixed effects, since review and permit-type composition may be endogenously determined by the regulatory environment.
The purpose of this stage is to remove systematic differences in permit timing attributable to observable ward characteristics so that the resulting strictness score isolates the alderman-specific component.

\begin{table}[htbp]
    \caption{Stage 1 Residualization: Permit-Level Log Processing Time on Ward Fundamentals}
    \label{tab:stage1_residualization}
    \input{../tasks/create_alderman_uncertainty_index/output/stage1_regression_ptfeTRUE_rtfeTRUE_porchTRUE_cafeFALSE_2stage_volLAG1_BOTH_through2022.tex}

    \footnotesize
    \textit{Notes:} Permit-level OLS regression from the baseline uncertainty-index specification, using permit applications through December 2022.
    Covariates are median household income, racial shares, homeownership, population, distance to the CBD, distance to the lake, CTA station density, and lag permit volume.
    Fixed effects are listed in the table.
    Regressions are unweighted.
    Standard errors in parentheses. * $p<0.10$, ** $p<0.05$, *** $p<0.01$.
\end{table}

\begin{figure}[htbp]
    \centering
    \includegraphics[width=0.8\textwidth]{../tasks/create_alderman_uncertainty_index/output/uncertainty_index_ptfeTRUE_rtfeTRUE_porchTRUE_cafeFALSE_2stage_volLAG1_BOTH_through2022.pdf}
    \caption{Distribution of Aldermanic Stringency Scores}
\figurenotes{Stringency scores are constructed from residualized permit processing times through December 2022 with empirical Bayes shrinkage.
Scores are standardized to mean zero and unit standard deviation.
Higher values indicate aldermen associated with longer permit processing times after controlling for ward characteristics and permit-level fixed effects.}
    \label{fig:aldermanic_strictness_scores}
\end{figure}

\subsection{Validation of the Strictness Index}
\label{app:strictness_validation}

The strictness index is intended to capture alderman-specific regulatory frictions rather than purely ward-level conditions.
Table~\ref{tab:permit_validation} shows that wards with higher stringency scores issue fewer high-discretion permits and exhibit more signs of project modification: a larger share of permits with any unit change and more units reduced per permit.
Figure~\ref{fig:predecessor_successor_strictness} complements that evidence by comparing predecessor and successor aldermen within the same ward.
The weak relationship indicates that the index is not simply reflecting time-invariant ward characteristics.

\begin{table}[htbp]
    \centering
    \caption{Validation of the Stringency Index Using Permit Outcomes}
    \label{tab:permit_validation}
    \input{../tasks/uncertainty_validation_checks/output/permit_validation_table_uncertainty_ptfeTRUE_rtfeTRUE_porchTRUE_cafeFALSE_2stage_volLAG1_BOTH_through2022.tex}

    \footnotesize
    \textit{Notes:} Ward-month regressions of permit outcomes on the standardized stringency index, using high-discretion permits through December 2022 and weights equal to ward-month permit counts.
    Outcomes are the number of high-discretion permits, the share of permits with any unit change, and units reduced per permit based on permit text.
    Standard errors in parentheses. * $p<0.10$, ** $p<0.05$, *** $p<0.01$.
\end{table}

\begin{figure}[htbp]
    \centering
    \includegraphics[width=0.72\textwidth]{../tasks/within_ward_strictness/output/predecessor_successor_scatter_uncertainty_ptfeTRUE_rtfeTRUE_porchTRUE_cafeFALSE_2stage_volLAG1_BOTH_through2022.pdf}
    \caption{Predecessor and Successor Alderman Strictness Within the Same Ward}
    \figurenotes{Each point compares the strictness score of an alderman to that of the predecessor or successor who represented the same ward in an adjacent period.
    The low correlation indicates that the strictness measure is not simply a persistent ward-level attribute.}
    \label{fig:predecessor_successor_strictness}
\end{figure}

\subsection{Placebo Checks}

Table~\ref{tab:low_discretion_placebo} asks whether the main stringency score predicts delays for low-discretion permits.
It does not.
Table~\ref{tab:low_discretion_density_placebo} carries the same idea into the density design by replacing the main score with a score built from residualized low-discretion processing times.
Those coefficients are small and statistically indistinguishable from zero.

\begin{table}[htbp]
    \centering
    \caption{Low-Discretion Processing Time as a Placebo Outcome}
    \label{tab:low_discretion_placebo}
    \input{../tasks/within_ward_strictness/output/low_discretion_falsification_compact_uncertainty_ptfeTRUE_rtfeTRUE_porchTRUE_cafeFALSE_2stage_volLAG1_BOTH_through2022.tex}
\end{table}

\begin{table}[htbp]
    \caption{Low-Discretion Score as a Placebo for Density Outcomes}
    \label{tab:low_discretion_density_placebo}

    \input{../tasks/new_construction_score_density_robustness/output/fe_table_500ft_all_zonegroup_segment_year_additive_clust_ward_pair_low_discretion_residualized_mainoutcomes.tex}

    \footnotesize
    \textit{Notes:} Same 500ft all-construction border-pair specification as Table~\ref{tab:density_main_table}, but replacing the main aldermanic stringency score with a placebo score built from alderman-level mean residualized low-discretion log processing times.
    Regressions include the same own-side controls, side-specific distance-to-boundary slopes, and fixed effects as the main density table.
    Regressions are unweighted.
    Standard errors clustered at the ward-pair level. * $p<0.10$, ** $p<0.05$, *** $p<0.01$.
\end{table}
